\subsection{Conditional Diffusion Model for Feature Synthesis}
\label{sec:diffusion}
At the core of the Diff-LM-GZSL framework is the Conditional Latent Diffusion Model (CLDM), which is specifically engineered to synthesize high-fidelity vibration features $x$ conditioned on the language-driven semantic embeddings $a$. Unlike traditional Generative Adversarial Networks (GANs) that often suffer from mode collapse and training instability in high-dimensional spectral spaces, diffusion models offer a more stable likelihood-based training objective and superior mode coverage. The generative process is modeled as a sequence of denoising steps that transform a simple Gaussian noise distribution into the complex manifold of mechanical fault features. By injecting the semantic vector $a$ derived from expert descriptions into this iterative process, the model learns to associate specific linguistic concepts—such as "periodic impulses" or "harmonic growth"—with their corresponding statistical patterns in the vibration feature domain.

The training of the CLDM involves two distinct phases: the forward diffusion process and the reverse denoising process. In the forward process, we gradually add Gaussian noise to a real feature vector $x_0$ sampled from the training set of seen classes. This process follows a Markov chain defined by a variance schedule $\{\beta_t\}_{t=1}^T$, where $T$ is the total number of diffusion steps. The state at time step $t$, denoted as $x_t$, can be expressed directly in terms of $x_0$ as $x_t = \sqrt{\bar{\alpha}_t} x_0 + \sqrt{1 - \bar{\alpha}_t} \epsilon$, where $\epsilon \sim \mathcal{N}(0, I)$ and $\bar{\alpha}_t = \prod_{i=1}^t (1 - \beta_i)$. As $t$ approaches $T$, the signal $x_T$ becomes approximately equivalent to pure isotropic Gaussian noise. The objective of the model is to learn the reverse transition $p_\theta(x_{t-1}|x_t, c)$, which effectively predicts the noise $\epsilon$ added at each step, given the current noisy state $x_t$ and a conditioning signal $c$. The training loss for the conditional diffusion model is defined as the mean squared error between the true noise and the predicted noise:
\begin{equation}
L_{diff} = E_{x, \epsilon \sim \mathcal{N}(0,I), t} \left[ ||\epsilon - \epsilon_\theta(x_t, t, c)||^2 \right]
\end{equation}
where $\epsilon_\theta$ is the denoising network parameterized by $\theta$. This loss function ensures that the model captures the reverse conditional density, allowing it to reconstruct features that are statistically indistinguishable from real measurements.

A critical component of the CLDM is the cross-modal condition injection mechanism, which bridges the gap between the linguistic semantic space and the vibration feature space. Since the semantic vector $a$ (from Eq. 6) and the latent feature $x$ reside in different latent manifolds, we introduce a conditioning transformer $\phi(\cdot)$ that maps the semantic embedding into a contextually relevant control signal $c$. We employ a cross-attention mechanism to perform this mapping, allowing the denoising network to selectively attend to specific dimensions of the semantic embedding during the reconstruction process. Formally, the conditioning signal $c$ is defined as:
\begin{equation}
c = \phi(a) = \text{Softmax}\left(\frac{Q_x K_a^T}{\sqrt{d_k}}\right) V_a
\end{equation}
where $Q_x$ represents the query derived from the noisy feature $x_t$, and $K_a, V_a$ are the keys and values projected from the semantic embedding $a$. Alternatively, for lower-dimensional feature spaces commonly encountered in 1D-CNN extractors, a Feature-wise Linear Modulation (FiLM) layer can be utilized, where $c$ acts as a scale and shift parameter applied to the intermediate activations of the denoising network. This explicit conditioning ensures that the generated features are not only realistic but also strictly aligned with the class-specific semantic description $T_y$.

The architecture of the denoising network $\epsilon_\theta$ is designed as a specialized U-Net variant or a deep Residual Network (ResNet) tailored for 1D feature vectors. Given the sequential nature of vibration features, we utilize a series of dilated convolutional blocks followed by cross-attention layers to capture both local dependencies and global semantic context. The network takes the noisy feature $x_t$, the time step $t$ (encoded via sinusoidal positional embeddings), and the conditioning signal $c$ as inputs. The specific architecture of the denoising network is defined as:
\begin{equation}
\epsilon_\theta(x_t, t, c) = \text{Decoder}\left( \text{Encoder}(x_t, \text{emb}(t)) \oplus c \right)
\end{equation}
where $\oplus$ facilitates the fusion of latent representations and conditioning context, and $\text{emb}(t)$ provides the temporal grounding for the denoising schedule. As illustrated in Fig. 4, the reverse diffusion process starts from $x_T \sim \mathcal{N}(0, I)$ and iteratively applies the denoising function $\epsilon_\theta$ guided by $c$ to produce a synthesized feature $x_{syn}$ after $T$ steps. This iterative refinement allows the model to "hallucinate" the missing spectral details for unseen faults by interpolating the learned relationships between language and vibration patterns. By training on a diverse set of seen classes, the CLDM develops a generalized understanding of how high-level mechanical descriptions map to low-level signal characteristics, which is the cornerstone of its zero-shot diagnostic capability.

\subsection{Semantic Alignment and Domain Shift Mitigation}
\label{sec:alignment}
A significant challenge in generalized zero-shot fault diagnosis is the "domain shift" or "distributional gap" between the features synthesized by the generative model and the real features extracted from physical sensors. Even with a powerful diffusion model, the synthesized features $x_{syn}$ for an unseen class may reside in a region of the feature space that is slightly offset from the actual (but unavailable during training) data distribution of that class. This shift can lead to the "seen-class bias" problem, where the final classifier tends to overfit the perfectly modeled seen classes while misclassifying unseen class samples. To mitigate this, we introduce a semantic alignment module that enforces structural and distributional consistency between the real and synthesized domains through a joint optimization of maximum mean discrepancy (MMD) and cycle-consistency constraints.

To ensure that the global distribution of generated features matches the underlying manifold of real vibration data, we utilize the Maximum Mean Discrepancy (MMD) as a non-parametric metric for domain alignment. MMD measures the distance between two distributions $P$ and $Q$ in a Reproducing Kernel Hilbert Space (RKHS) $\mathcal{H}$. For the diagnosed vibration features, let $X_{real} = \{f_{real}^{(i)}\}$ be the set of real features from seen classes and $X_{syn} = \{f_{syn}^{(j)}\}$ be the set of features synthesized by the CLDM for the same classes during the alignment phase. The MMD objective is formulated as:
\begin{equation}
\text{MMD}(P, Q) = \left\| \frac{1}{n} \sum_{i=1}^n \psi(f_{real}^{(i)}) - \frac{1}{m} \sum_{j=1}^m \psi(f_{syn}^{(j)}) \right\|_{\mathcal{H}}
\end{equation}
where $\psi(\cdot)$ is a feature mapping to the RKHS, typically implemented using a multi-scale Gaussian RBF kernel. By minimizing this term, we force the diffusion model to generate features that are statistically consistent with the training data, thereby narrowing the gap between "hallucinated" and "observed" fault signatures.

The semantic alignment is further reinforced by a classification-aware loss that ensures the synthesized features preserve the categorical boundaries defined in the semantic space. We introduce an alignment loss $L_{align}$ which combines the MMD metric with a supervised classification term on the synthesized data:
\begin{equation}
L_{align} = \text{MMD}(f_{real}, f_{syn}) + \lambda L_{cls}(f_{syn}, y)
\end{equation}
where $L_{cls}$ is the cross-entropy loss and $\lambda$ is a balancing coefficient. This ensures that the generated features for class $y$ are not only located in the correct region of the feature space but are also discriminative enough to be correctly identified by a classifier trained on real features. Furthermore, to prevent the loss of semantic information during the diffusion-denoising process, we implement a cycle-consistency or reconstruction loss. This loss requires that a feature $x$, once transformed into the latent space and reconstructed, should be recoverable:
\begin{equation}
L_{recon} = \|x - G(E(x))\|^2
\end{equation}
where $E$ and $G$ represent the encoding and generative components of the CLDM pipeline. This constraint maintains the structural integrity of the feature manifold across multiple diffusion steps.

The total training objective for the Diff-LM-GZSL framework integrates the diffusion likelihood, the domain alignment, and the reconstruction constraints into a single unified loss function:
\begin{equation}
L_{total} = L_{diff} + \alpha L_{align} + \beta L_{recon}
\end{equation}
where $\alpha$ and $\beta$ are hyper-parameters that control the influence of domain shift mitigation and structural preservation, respectively. The role of $\alpha$ is particularly crucial; as shown in the domain shift visualization in Fig. 10, a higher $\alpha$ forces the synthesized unseen class distributions to align more closely with the shared semantic manifold of the seen classes, effectively regularizing the model against overfitting to specific training set noise. By optimizing $L_{total}$ across all seen categories, the framework learns a robust, semantically-grounded mapping that can generalize to entirely new fault descriptions without catastrophic forgetting or distributional collapse. This holistic alignment ensures that when the model encounters a description of a "bearing cage failure" for the first time, its synthesized features will occupy the correct neighborhood relative to known "ball" or "race" failures.

\subsection{Final Classification and Inference}
\label{sec:inference}
Once the conditional latent diffusion model and the alignment module have been fully optimized on the seen fault classes, the framework is ready to perform zero-shot diagnosis. The inference process begins with the generation of a representative set of synthetic features for the unseen classes. For each unseen fault category $u \in \mathcal{Y}_{unseen}$, we take its standardized linguistic description $T_u$ and project it into the semantic space to obtain the embedding $a_u$. This embedding is then used as the conditioning signal $c_u$ for the CLDM. By sampling $N$ different noise vectors from $\mathcal{N}(0, I)$ and passing them through the reverse diffusion process $T$ times, we synthesize a comprehensive set of "pseudo-features" $\mathcal{X}_{syn, unseen}$ that capture the expected statistical variance of the unseen fault. This process effectively expands the training set from the seen label space $\mathcal{Y}_{seen}$ to the generalized label space $\mathcal{Y}_{seen} \cup \mathcal{Y}_{unseen}$.

In the generalized zero-shot learning (GZSL) setting, the model must be able to classify a test sample $x_{test}$ into any category within the union of seen and unseen classes. To achieve this, we construct a global feature matrix $F$ by concatenating the real features from seen classes and the synthesized features for both seen and unseen classes:
\begin{equation}
F = [f_{seen}^{(1)}, \dots, f_{seen}^{(M)}; f_{syn, unseen}^{(1)}, \dots, f_{syn, unseen}^{(N)}]
\end{equation}
where $M$ and $N$ represent the number of samples for seen and unseen classes, respectively. A final linear classifier with weights $W_c$ and bias $b_c$ is then trained on this augmented dataset. During the testing phase, the predicted health state $y^*$ for an input vibration feature $f_{test}$ is determined using the maximum softmax probability across the entire label space:
\begin{equation}
y^* = \text{argmax } \text{softmax}(W_c f_{test} + b_c)
\end{equation}
This formulation ensures that the model can handle the competition between seen and unseen classes fairly. The effectiveness of this synthesized-feature-driven classification is quantitatively reflected in Table X, where the MMD scores indicate a high degree of overlap between actual unseen distributions and our generated pseudo-features.

The complete operational logic of the Diff-LM-GZSL framework is summarized in Algorithm 1. The training phase focuses on the joint optimization of the feature extractor, the LLM-based semantic projector, and the CLDM, while the inference phase centers on the stochastic generation and classification of unseen signatures. The time complexity of the training phase is primarily dominated by the $T$ steps of the diffusion process per sample, while the inference phase's complexity for a single real-time diagnosis is extremely low, as it only requires a single pass through the 1D-CNN feature extractor and the final softmax classifier. The synthetic feature generation is performed offline or during a one-time "adaptation" phase whenever a new fault description is added to the system.

\begin{algorithm}
\caption{Diff-LM-GZSL: Training and Inference}
\begin{algorithmic}[1]
\State \textbf{Input:} Raw signals $X_{seen}$ with labels $Y_{seen}$, text descriptions $\{T_y\}$ for all $y \in \mathcal{Y}_{seen} \cup \mathcal{Y}_{unseen}$
\State \textbf{Training (on Seen Classes):}
\State \quad Extract features $x = f_\phi(X_{seen})$ via 1D-CNN (Eq. 4)
\State \quad Compute semantic embeddings $a = W_s \text{Encoder}(T) + b_s$ (Eq. 5-6)
\State \quad \textbf{repeat}
\State \quad \quad Sample $x_0 \sim q(x)$, $t \sim \text{Uniform}(1, \dots, T)$, $\epsilon \sim \mathcal{N}(0, I)$
\State \quad \quad Update $\theta$ by minimizing $L_{total}$ (Eq. 13) including $L_{diff}$, $L_{align}$, $L_{recon}$
\State \quad \textbf{until} converged
\State \textbf{Zero-Shot Synthesis (for Unseen Classes):}
\State \quad \textbf{for} each $y \in \mathcal{Y}_{unseen}$ \textbf{do}
\State \quad \quad Compute $c = \phi(a_y)$ using cross-attention (Eq. 8)
\State \quad \quad Generate $N$ features $x_{syn, y}$ via reverse diffusion $p_\theta(x_{t-1}|x_t, c)$
\State \quad \textbf{end for}
\State \textbf{Inference:}
\State \quad Train a final softmax classifier on $F = [X_{seen}; X_{syn, unseen}]$ (Eq. 14)
\State \quad \textbf{Predict:} $y^* = \text{argmax } \text{softmax}(W_c f_{test} + b_c)$ (Eq. 15)
\end{algorithmic}
\end{algorithm}

This multi-stage pipeline provides a robust solution for the evolving nature of industrial assets, where new failure modes emerge without historical precedents. By leveraging the vast linguistic knowledge of LLMs and the generative precision of latent diffusion models, Diff-LM-GZSL bridges the semantic-physical divide, enabling continuous and accurate fault diagnosis in zero-shot scenarios. The resulting model is not only capable of identifying unseen faults but also providing a technical justification through its reliance on expert language descriptions.
